\documentclass[a4,11pt]{aleph-notas}

% -- Paquetes adicionales
\usepackage{aleph-comandos}
\usepackage{booktabs}
\usepackage{pdflscape}

% -- Datos   
\institucion{Facultad de Ciencias Exactas, Naturales y Ambientales}
\carrera{Catálogo STEM}
\asignatura{Matemática III}
\tema{Video exposición no. 2: Espacios vectoriales y bases}
\autor{Andrés Merino}
\fecha{Periodo 2026-1}

\logouno[0.16\textwidth]{Logos/logoPUCE_04_ac}
\definecolor{colortext}{HTML}{1957d1}
\definecolor{colordef}{HTML}{1957d1}
\fuente{montserrat}


% -- Otros comandos




\begin{document}

\encabezado


%%%%%%%%%%%%%%%%%%%%%%%%%%%%%%%%%%%%%%%%
\section{Resultado de aprendizaje}
%%%%%%%%%%%%%%%%%%%%%%%%%%%%%%%%%%%%%%%%

\begin{itemize}
\item 
    Esta actividad evalúa si el estudiante ({Criterio 1.3:})  Comprende los conceptos de espacios vectoriales y aplicaciones lineales, incluyendo la base y dimensión, transformaciones lineales y sus propiedades.
\end{itemize}

%%%%%%%%%%%%%%%%%%%%%%%%%%%%%%%%%%%%%%%%
\section{Indicaciones}
%%%%%%%%%%%%%%%%%%%%%%%%%%%%%%%%%%%%%%%%

\begin{itemize}
\item 
    El estudiante deberá elaborar diapositivas profesionales o un apoyo visual equivalente y grabar un video presentándolas, con el rostro visible durante toda la exposición.
\item 
    No se aceptará un video en el que el estudiante únicamente lea las diapositivas; se deberá evidenciar dominio del tema mediante una explicación clara y personal.
\item 
    No se aceptarán presentaciones ni libretos claramente generado por inteligencia artificial y descontextualizado de lo visto en clase.
\item 
    La duración máxima del video será de 8 minutos.
\end{itemize}


%%%%%%%%%%%%%%%%%%%%%%%%%%%%%%%%%%%%%%%%
\section{Descripción}
%%%%%%%%%%%%%%%%%%%%%%%%%%%%%%%%%%%%%%%%

\begin{enumerate}
    \item {Elaboración de diapositivas o apoyo visual:}
    \begin{itemize}
        \item Crear una presentación clara y ordenada sobre espacios vectoriales, subespacios, base y dimensión, puede ser escrito a mano siempre que sea legible.
        \item Incluir definiciones y propiedades clave: subespacio vectorial, base, dimensión.
        \item Presentar de forma explícita los tres ejercicios solicitados y su resolución paso a paso.
        \item El apoyo visual debe mostrar operaciones y justificaciones, no solo respuestas finales.
    \end{itemize}

    \item {Ejercicios a resolver en el video:}
    \begin{itemize}
        \item \textbf{Ejercicio 1:} En $\mathbb{R}^{2\times 2}$, sea
        \[
        W=\left\{\begin{pmatrix}a&b\\c&d\end{pmatrix}\in\mathbb{R}^{2\times2}:\ a+d-2c=1,\ \ a+2b-c=0\right\}.
        \]
        Determinar si $W$ es subespacio vectorial de $\mathbb{R}^{2\times2}$.

        \item \textbf{Ejercicio 2:} En $\mathbb{R}^{2\times 2}$, sea
        \[
        F=\left\{\begin{pmatrix}a&b\\c&d\end{pmatrix}\in\mathbb{R}^{2\times2}:\ a+3b+d=0,\ \ a+2b-c=0\right\}.
        \]
            Verificar que $F$ es subespacio vectorial de $\mathbb{R}^{2\times2}$.
        \item \textbf{Ejercicio 3:}  Determinar una base de $F$ e indicar su dimensión.
    \end{itemize}

    \item {Guía de explicación paso a paso:}
    \begin{itemize}
        \item En cada ejercicio, enunciar el conjunto y el objetivo antes de empezar a operar.
        \item Mostrar cada paso algebraico relevante y explicar por qué ese paso permite concluir si es o no subespacio.
        \item En el Ejercicio 3, indicar cómo se construyen los generadores, cómo se verifica independencia lineal y cómo se concluye la dimensión.
    \end{itemize}

    \item {Grabación del video:}
    \begin{itemize}
        \item El estudiante deberá presentarse en cámara durante toda la exposición.
        \item El rostro debe ser visible en todo momento.
        \item No se permite únicamente voz en off ni una grabación en la que solo se lean las diapositivas.
    \end{itemize}

    \item {Estructura mínima:}
    \begin{itemize}
        \item {Introducción:} explicación breve sobre espacios vectoriales, subespacios y la importancia de base y dimensión.
        \item {Desarrollo:} resolución completa de los tres ejercicios con explicación paso a paso, usando el apoyo visual.
        \item {Cierre:} síntesis final sobre los criterios para determinar subespacios y para hallar base y dimensión.
    \end{itemize}
\end{enumerate}

%%%%%%%%%%%%%%%%%%%%%%%%%%%%%%%%%%%%%%%%
\section{Entrega}
%%%%%%%%%%%%%%%%%%%%%%%%%%%%%%%%%%%%%%%%

Se deberá entregar en el aula virtual el enlace de \textbf{OneDrive} del video con permisos de acceso habilitados. Si el enlace no permite visualizar el material, la actividad no podrá ser calificada y se asignará \textbf{0}.

%%%%%%%%%%%%%%%%%%%%%%%%%%%%%%%%%%%%%%%%
\section{Rúbrica de calificación}
%%%%%%%%%%%%%%%%%%%%%%%%%%%%%%%%%%%%%%%%

En caso de presentar un video con voz en off, rostro no visible o fuera del rango de tiempo establecido, se reducirá un nivel a cada punto de la rúbrica.

\begin{landscape}
\begin{center}\footnotesize
\begin{tabular}{p{5cm}p{4cm}p{4cm}p{4cm}p{4cm}}
\toprule
\textbf{Indicador} & \textbf{Excelente} & \textbf{Bueno} & \textbf{Aceptable} & \textbf{Insuficiente} \\ \midrule
\textbf{Ejercicio 1: $W$ no es subespacio (10 pts)} & Presenta correctamente $W$ y demuestra con contraejemplo de suma o cierre escalar que no es subespacio (10) & Resolución correcta con una omisión menor en la justificación (8) & Identifica que no es subespacio, pero la demostración es parcial (6) & No demuestra correctamente o concluye de forma errónea (0) \\ \midrule
\textbf{Ejercicio 2: $F$ es subespacio (15 pts)} & Verifica correctamente que $F$ es subespacio a partir de las ecuaciones dadas, con justificación completa (15) & Desarrollo correcto con detalles menores faltantes en la justificación (12) & Verificación parcialmente correcta o con pasos incompletos (9) & Errores conceptuales importantes o sin justificación suficiente (0) \\ \midrule
\textbf{Ejercicio 3: base y dimensión de $F$ (15 pts)} & Determina una base válida de $F$ e indica explícitamente su dimensión con justificación completa (15) & Desarrollo correcto con detalles menores faltantes en la justificación (12) & Base o dimensión parcialmente correctas (9) & Errores conceptuales importantes o sin justificación suficiente (0) \\ \midrule
\textbf{Presentación en video (5 pts)}
& Exposición clara y fluida; rostro visible durante toda la presentación (5)
& Exposición clara con ligeras interrupciones; rostro visible (4)
& Exposición poco fluida o lectura excesiva; rostro visible parcial (3)
& Voz en off o rostro no visible / exposición ininteligible (0) \\ \midrule
\textbf{Gestión del tiempo (5 pts)} & 8:00 $\pm$ 30 s (5) & 8:00 $\pm$ 60 s (4) & 8:00 $\pm$ 90 s (3) & Fuera de $\pm$ 90 s (0) \\ \bottomrule
\end{tabular}
\end{center}
\end{landscape}



\end{document}
